\documentclass[12pt,leqno]{article}

\usepackage{amsmath}
\usepackage{amssymb}
\usepackage{bm}
\usepackage{enumitem}
\usepackage{hyperref}
\usepackage{logicproof}
\usepackage{synttree}
\usepackage{xcolor}

\setlength{\textheight}{9.1in}
\setlength{\topmargin}{-.1in}
\setlength{\headsep}{0in}
\setlength{\oddsidemargin}{-.1in}
\setlength{\textwidth}{6.5in}

\setlength{\parskip}{1em}
\setlength{\parindent}{0pt}

% for logicproof
\setlength{\subproofhorizspace}{.5em}
\setlength{\intersubproofvertspace}{.5em}

% custom symbols
\newcommand{\DefEq}{\stackrel{\text{def}}{=}}

% proof rules
\newcommand{\Intro}[1]{{#1}{\text{I}}}
\newcommand{\IntroA}[1]{{#1}{\text{I}_1}}
\newcommand{\IntroB}[1]{{#1}{\text{I}_2}}
\newcommand{\Elim}[1]{{#1}{\text{E}}}
\newcommand{\ElimA}[1]{{#1}{\text{E}_1}}
\newcommand{\ElimB}[1]{{#1}{\text{E}_2}}

% truth tables
\newcommand{\T}{\mathtt{T}}
\newcommand{\F}{\mathtt{F}}

\title{CS 511 Homework \#3}
\author{Anthony DeRossi}
\date{26 Sep 2024}

\begin{document}
\maketitle

\section*{Exercise 1}

Go to page 13 in \textbf{\textit{Lecture Slides 09}}. Your task is to carefully
do part 1 of the exercise on that page.

Write a natural-deduction proof of the following WFF:

\[
  \varphi_1 \triangleq \lnot (p \land q \land r) \to
    (\lnot p \lor \lnot q \lor \lnot r)
\]

\newpage
\begin{logicproof}{3}
  \begin{subproof}
    \lnot ((p \land q) \land r)                   & assumption \\
    \begin{subproof}
      \lnot ((\lnot p \lor \lnot q) \lor \lnot r) & assumption \\
      \begin{subproof}
        \lnot p                                   & assumption \\
        \lnot p \lor \lnot q                      & $\IntroB{\lor}$ 3 \\
        (\lnot p \lor \lnot q) \lor \lnot r       & $\IntroB{\lor}$ 4 \\
        \bot                                      & $\Elim{\lnot}$ 2, 5
      \end{subproof}
      \lnot \lnot p                               & $\Intro{\lnot}$ 3-6 \\
      \begin{subproof}
        \lnot q                                   & assumption \\
        \lnot p \lor \lnot q                      & $\IntroA{\lor}$ 8 \\
        (\lnot p \lor \lnot q) \lor \lnot r       & $\IntroB{\lor}$ 9 \\
        \bot                                      & $\Elim{\lnot}$ 2, 10
      \end{subproof}
      \lnot \lnot q                               & $\Intro{\lnot}$ 8-11 \\
      \begin{subproof}
        \lnot r                                   & assumption \\
        (\lnot p \lor \lnot q) \lor \lnot r       & $\IntroA{\lor}$ 13 \\
        \bot                                      & $\Elim{\lnot}$ 2, 14
      \end{subproof}
      \lnot \lnot r                               & $\Intro{\lnot}$ 13-15 \\
      p                                           & $\Elim{\lnot\lnot}$ 7 \\
      q                                           & $\Elim{\lnot\lnot}$ 12 \\
      r                                           & $\Elim{\lnot\lnot}$ 16 \\
      p \land q                                   & $\Intro{\land}$ 17, 18 \\
      (p \land q) \land r                         & $\Intro{\land}$ 20, 19 \\
      \bot                                        & $\Elim{\lnot}$ 1, 21
    \end{subproof}
    \lnot \lnot ((\lnot p \lor \lnot q) \lor \lnot r)
      & $\Intro{\lnot}$ 2-22 \\
    (\lnot p \lor \lnot q) \lor \lnot r           & $\Elim{\lnot\lnot}$ 23
  \end{subproof}
  \lnot ((p \land q) \land r) \to ((\lnot p \lor \lnot q) \lor \lnot r)
    & $\Intro{\to}$ 1-24
\end{logicproof}

\newpage
\section*{Exercise 2}

Go to page 18 in \textbf{\textit{Lecture Slides 10}}. Your task is to carefully
do part 1 of the exercise on that page.

Use the tableaux method to show the validity of the following more general
version of de Morgan's law:

\[
  \varphi_1 \triangleq \lnot (p \land q \land r) \to
    (\lnot p \lor \lnot q \lor \lnot r)
\]

\medskip
We show the formula is valid by showing that its negation is unsatisfiable:

\begin{center}
  \synttree{8 \branchheight{.3in} \childsidesep{5em} \childattachsep{1em}}
  [ $\lnot (\lnot ((p \land q) \land r) \to
      ((\lnot p \lor \lnot q) \lor \lnot r))$
    [ $\lnot ((p \land q) \land r)$
      [ $\lnot ((\lnot p \lor \lnot q) \lor \lnot r)$
        [ $\lnot (\lnot p \lor \lnot q)$
          [ $\lnot \lnot r$
            [ $\lnot \lnot p$
              [ $\lnot \lnot q$
                [ $\color{blue} \bm r$
                  [ $\color{green} \bm p$
                    [ $\color{red} \bm q$
                      [ $\lnot (p \land q)$
                        [ $\color{green} \bm{\lnot p}$
                          [ \sf X ]
                        ]
                        [ $\color{red} \bm{\lnot q}$
                          [ \sf X ]
                        ]
                      ]
                      [ $\color{blue} \bm{\lnot r}$
                        [ \sf X ]
                      ]
                    ]
                  ]
                ]
              ]
            ]
          ]
        ]
      ]
    ]
  ]
\end{center}

All paths are closed, therefore this formula is unsatisfiable and the original
formula is valid.

\newpage
\section*{Problem 1}

\begin{itemize}
\item[(a)] Go to page 13 in \textbf{\textit{Lecture Slides 09}} once more.
  Your task is to carefully do parts 2, 3, and 4 of the exercise on that page.
\item[(b)] Go to page 18 in \textbf{\textit{Lecture Slides 10}} once more.
  Your task is to carefully do parts 2, 3, and 4 of the exercise on that page.
\end{itemize}

% TODO

\vspace{3\baselineskip}
\hrule

Solutions to the Lean 4 exercises are available on GitHub:

\url{https://github.com/anthonyde/math2001/blob/cs511-fall2024/Math2001/CS511_Homework/hw3.lean}

\section*{Exercise 3}

\section*{Exercise 4}

\section*{Problem 2}

\end{document}
